% 02_agents.tex - Robot Agents Module
% AlohaMini Codebase Documentation

\section{Agents Module}

% ============================================
% SLIDE: Module Overview
% ============================================
\begin{frame}{agents/aloha\_mini/ Module}
\begin{columns}
\begin{column}{0.45\textwidth}
    \textbf{Purpose:} Define robot agents for ManiSkill3 simulation

    \vspace{0.3cm}

    \textbf{Files:}
    \begin{itemize}
        \item \filepath{base\_agent.py} (208 lines)
        \item \filepath{aloha\_mini\_so100\_v2.py} (351 lines)
        \item \filepath{\_\_init\_\_.py}
    \end{itemize}
\end{column}
\begin{column}{0.55\textwidth}
    \centering
    \begin{tikzpicture}[scale=0.75, transform shape]
        % ManiSkill BaseAgent (top)
        \node[submodule, text width=3cm] (maniskill) at (0,3) {\scriptsize mani\_skill.BaseAgent};

        % AlohaMiniBase (middle)
        \node[class, text width=3cm] (base) at (0,1.5) {
            \textbf{AlohaMiniBase}\\
            \scriptsize (Abstract)
        };

        % SO100V2 (bottom)
        \node[class, text width=2.8cm] (so100) at (0,0) {
            SO100V2\\
            \scriptsize (16 DOF)
        };

        % Inheritance arrows - directly connected
        \draw[arrow] (base) -- (maniskill);
        \draw[arrow] (so100) -- (base);
    \end{tikzpicture}
\end{column}
\end{columns}
\end{frame}

% ============================================
% SLIDE: base_agent.py Overview
% ============================================
\begin{frame}[fragile]{base\_agent.py - Base Class}
\begin{columns}
\begin{column}{0.5\textwidth}
\begin{lstlisting}[basicstyle=\ttfamily\scriptsize]
class AlohaMiniBaseAgent(BaseAgent):
    """Base agent for variants."""
    BASE_JOINT_NAMES = [
        "root_x_axis_joint",
        "root_y_axis_joint",
        "root_z_rotation_joint",
    ]
    LIFT_JOINT_NAMES = ["vertical_move"]

    @abstractmethod
    def get_left_ee_pose(self): pass

    @abstractmethod
    def get_right_ee_pose(self): pass
\end{lstlisting}
\end{column}
\begin{column}{0.5\textwidth}
    \textbf{Common Functionality:}
    \begin{itemize}
        \item Joint name definitions
        \item Controller parameter defaults
        \item \code{\_create\_base\_controller()}
        \item \code{is\_grasping()} (Template Method)
        \item \code{is\_static()}
    \end{itemize}

    \vspace{0.2cm}

    \begin{block}{Design Pattern}
    \footnotesize\textbf{Template Method}: \code{is\_grasping()} calls abstract \code{\_check\_single\_arm\_grasping()}
    \end{block}
\end{column}
\end{columns}
\end{frame}

% ============================================
% SLIDE: Controller Helper Methods
% ============================================
\begin{frame}[fragile]{base\_agent.py - Controller Helpers}
\begin{lstlisting}
def _create_base_controller(self):
    """Create PDBaseVelController for virtual mobile base."""
    return PDBaseVelControllerConfig(
        self.base_joint_names,
        lower=[-1, -1, -3.14],    # [vx, vy, omega]
        upper=[1, 1, 3.14],
        damping=self.base_damping,        # 1000
        force_limit=self.base_force_limit, # 500
    )

def _create_lift_pos_controller(self):
    """Create position controller for lift mechanism."""
    return PDJointPosControllerConfig(
        self.lift_joint_names,
        stiffness=self.lift_stiffness,    # 2e3
        damping=self.lift_damping,        # 2e2
        force_limit=self.lift_force_limit, # 150
        normalize_action=False,
    )
\end{lstlisting}
\end{frame}

% ============================================
% SLIDE: is_grasping Template Method
% ============================================
\begin{frame}[fragile]{base\_agent.py - Grasping Detection}
\begin{lstlisting}
def is_grasping(self, object: Actor, min_force=0.5,
                max_angle=110, arm_id=None):
    """Check if robot is grasping an object."""
    if arm_id is None:
        # Check both arms
        arm1 = self._check_single_arm_grasping(
            object, min_force, max_angle, arm_id=1)
        arm2 = self._check_single_arm_grasping(
            object, min_force, max_angle, arm_id=2)
        return torch.logical_or(arm1, arm2)
    else:
        return self._check_single_arm_grasping(
            object, min_force, max_angle, arm_id)

@abstractmethod
def _check_single_arm_grasping(self, object, ...):
    """Implemented by subclass (different gripper structures)"""
    pass
\end{lstlisting}

\begin{block}{Note}
Different gripper structures require different contact force calculations.
\end{block}
\end{frame}

% ============================================
% SLIDE: aloha_mini_so100_v2.py Overview
% ============================================
\begin{frame}{aloha\_mini\_so100\_v2.py - SO100 Variant}
\begin{columns}
\begin{column}{0.5\textwidth}
    \textbf{Configuration:}
    \begin{itemize}
        \item \code{uid = "aloha\_mini\_so100\_v2"}
        \item Virtual base (prismatic X/Y)
        \item 16 DOF total
        \item SO100 arm design
    \end{itemize}

    \vspace{0.3cm}

    \textbf{Joint Names:}
    \begin{itemize}
        \item \code{left\_shoulder\_pan/lift}
        \item \code{elbow\_flex}, \code{wrist\_flex}
        \item \code{wrist\_roll}, \code{gripper}
    \end{itemize}
\end{column}
\begin{column}{0.5\textwidth}
    \textbf{Controller Parameters:}
    \begin{table}
    \centering
    \footnotesize
    \begin{tabular}{lr}
    \toprule
    \textbf{Parameter} & \textbf{Value} \\
    \midrule
    arm\_stiffness & 3e3 \\
    arm\_damping & 2e2 \\
    arm\_force\_limit & 100 N \\
    gripper\_stiffness & 100 \\
    \bottomrule
    \end{tabular}
    \end{table}

    \vspace{0.3cm}

    \textbf{Cameras:}
    \begin{itemize}
        \item cam\_main (320$\times$240)
        \item cam\_left\_wrist (128$\times$128)
        \item cam\_right\_wrist (128$\times$128)
    \end{itemize}
\end{column}
\end{columns}
\end{frame}

% ============================================
% SLIDE: Keyframes
% ============================================
\begin{frame}[fragile]{Keyframes Configuration}
\begin{lstlisting}[basicstyle=\ttfamily\scriptsize]
keyframes = dict(
    rest=Keyframe(
        qpos=np.array([
            0.0, 0.0, 0.0,  # Base: x, y, rotation
            0.0,            # Lift
            0.0, 0.0, 0.0, 0.0, 0.0, 0.0,  # Left arm
            0.0, 0.0, 0.0, 0.0, 0.0, 0.0,  # Right arm
        ]),
        pose=sapien.Pose(p=[0, 0, 0]),
    ),
    ready=Keyframe(
        qpos=np.array([
            0.0, 0.0, 0.0, 0.05,           # base + lift
            0.0, 0.3, -0.3, 0.0, 0.3, 0.0, # left
            0.0, 0.3, -0.3, 0.0, 0.3, 0.0, # right
        ]),
    ),
)
\end{lstlisting}
\end{frame}

% ============================================
% SLIDE: Camera Configuration
% ============================================
\begin{frame}[fragile]{Camera Setup}
\begin{lstlisting}[basicstyle=\ttfamily\scriptsize]
@property
def _sensor_configs(self):
    q_main = euler_to_quat_xyz(
        math.radians(90), math.radians(90), 0)  # overhead
    q_wrist = euler_to_quat_xyz(
        0, math.radians(45), math.radians(-90)) # 45deg down

    return [
        CameraConfig(uid="cam_main",
            pose=Pose.create_from_pq(p=[0.0, 0.0, 1.45], q=q_main),
            width=320, height=240, fov=1.5, entity_uid="base_link"),
        CameraConfig(uid="cam_left_wrist",
            pose=Pose.create_from_pq(p=[0.0, 0.0, 0.15], q=q_wrist),
            width=128, height=128, fov=1.3, entity_uid="left_link4"),
    ]
\end{lstlisting}
\end{frame}

% ============================================
% SLIDE: SO100 Grasping Detection
% ============================================
\begin{frame}[fragile]{Grasping Detection}
\begin{lstlisting}
def _check_single_arm_grasping(self, object, min_force=0.5,
                                max_angle=110, arm_id=1):
    """SO100 style: two-finger grasping detection."""
    finger1_link = self.left_finger1_link  # Fixed_Jaw
    finger2_link = self.left_finger2_link  # Moving_Jaw

    l_contact = self.scene.get_pairwise_contact_forces(
        finger1_link, object)
    r_contact = self.scene.get_pairwise_contact_forces(
        finger2_link, object)

    lforce = torch.linalg.norm(l_contact, axis=1)
    rforce = torch.linalg.norm(r_contact, axis=1)

    # Direction to open gripper
    ldirection = finger1_link.pose.to_transformation_matrix()[..., :3, 1]
    rdirection = -finger2_link.pose.to_transformation_matrix()[..., :3, 1]

    # Check force and angle thresholds
    lflag = (lforce >= min_force) & (angle <= max_angle)
    rflag = (rforce >= min_force) & (angle <= max_angle)
    return torch.logical_and(lflag, rflag)
\end{lstlisting}
\end{frame}

% ============================================
% SLIDE: Class Diagram Summary
% ============================================
\begin{frame}{Agents Module - Class Diagram}
\centering
\begin{tikzpicture}[scale=0.7, transform shape]
    % BaseAgent from ManiSkill
    \node[class, text width=3cm, fill=green!10] (base) at (0,5) {
        \footnotesize\textbf{BaseAgent}\\
        \tiny (mani\_skill)
    };

    % AlohaMiniBaseAgent
    \node[class, text width=3.5cm, minimum height=2cm] (aloha) at (0,2.2) {
        \footnotesize\textbf{AlohaMiniBaseAgent}\\
        \hrule
        \tiny BASE\_JOINT\_NAMES\\
        \tiny LIFT\_JOINT\_NAMES\\
        \tiny is\_grasping()\\
        \tiny @abstract get\_ee\_pose()
    };

    % AlohaMiniSO100V2
    \node[class, text width=3.2cm, minimum height=1.6cm] (so100) at (0,-0.8) {
        \footnotesize\textbf{AlohaMiniSO100V2}\\
        \tiny uid = "aloha\_mini\_so100\_v2"\\
        \hrule
        \tiny 16 DOF (base+lift+arms)\\
        \tiny \_check\_single\_arm\_grasping()
    };

    % Arrows
    \draw[arrow] (aloha) -- (base);
    \draw[arrow] (so100) -- (aloha);
\end{tikzpicture}
\end{frame}
